\chapter{Relación del Fenómeno con el Ser Humano}
\label{sec.cap3:ser_humano}

Las consecuencias de un incendio industrial sobre las personas abarcan desde traumatismos y quemaduras térmicas cutáneas
hasta la asfixia celular derivada de los subproductos de la pirólisis. En la práctica pericial y forense se constata que
la inmensa mayoría de las víctimas mortales en incendios estructurales no fallece por la acción directa del fuego o de
las llamas, sino por la inhalación de atmósferas contaminadas con gases tóxicos y partículas en suspensión.

\section{Toxicidad del Monóxido de Carbono y el Cianuro}

Bajo condiciones reales de incendio en recintos cerrados o semi-confinados, la combustión se desarrolla de manera
incompleta debido a la restricción progresiva en el aporte de oxígeno. Esta deficiencia, combinada con la degradación
térmica de los materiales, se liberan concentraciones letales de gases asfixiantes, entre los cuales destacan el
monóxido de carbono y el cianuro de hidrógeno:

\begin{itemize}
    \item Monóxido de Carbono (CO): Es un gas inodoro, incoloro e insípido, lo que impide su percepción sensorial directa
          por los ocupantes. Al ingresar por el tracto respiratorio, atraviesa los alvéolos pulmonares y se enlaza con el
          hierro de la hemoglobina sanguínea con una afinidad aproximadamente $250$ veces superior a la del oxígeno
          molecular. Este fenómeno desplaza al oxígeno y bloquea su transporte. La disminucion progresiva del oxigeno
          produce alteraciones psicomotrices, perdida del conocimiento, coma y colapsos cardiovasculares.
    \item Cianuro de Hidrógeno (HCN): Se desprende durante la descomposición pirolítica de materiales que contienen
          átomos de nitrógeno en su estructura molecular (lanas, poliuretano, nylon y espumas de aislación). El gas
          cianuro es rápidamente absorbido y actúa a nivel celular como un potente inhibidor enzimático, volviendo
          incapaces a las celulas de usar el oxigeno en sangre, produciendo muerte celular.
\end{itemize}

La inhalación combinada de monóxido de carbono y cianuro ejerce un efecto toxicológico sinérgico y multiplicativo. La
primera manifestación clínica consiste en la pérdida de la agudeza visual, confusión mental y desorientación espacial
aguda. Este cuadro anula la capacidad del individuo para orientarse o accionar mecanismos de escape, provocando debilidad
muscular súbita en miembros inferiores que postra a la víctima en el plano del suelo antes de poder alcanzar una salida de
emergencia.

\section{Dinámica del Humo, Pérdida de Visibilidad y Comportamiento Humano}

Además de la asfixia química inducida por monóxido de carbono y cianuro de hidrógeno, el humo de combustión ejerce
mecanismos físicos y psicológicos que degradan la capacidad de supervivencia de los ocupantes:

\begin{itemize}
    \item Gases Irritantes y Corrosivos: La descomposición de materiales presentes en aislaciones de cables libera
          compuestos que atacan de inmediato las mucosas oculares y las vías respiratorias, induciendo lagrimeo, cierre
          involuntario de los párpados y espasmos que impiden mantener los ojos abiertos o respirar a frecuencias
          normales.
    \item Temperatura de la Mezcla Gaseosa: La inhalación de gases a temperaturas superiores a $120\Celsius$ provoca
          quemaduras en los tractos respiratorios, causando la muerte por colapso respiratorio mecánico.
    \item Oscurecimiento Óptico y Velocidad de Marcha: Las partículas de hollín suspendidas disminuyen la luz
          ambiental y la señalización de emergencia, desorientando al individuo y obligandolo a reducir su marcha para
          poder ubicarse.
\end{itemize}

El entendimiento riguroso de estas limitaciones fundamenta la obligatoriedad de los sistemas de evacuación y
compartimentación exigidos por la legislación laboral, así como la adopción de tecnologías de detección ultraprecoz y
agentes extintores limpios desarrollados en los capítulos subsiguientes.

